\documentclass[11pt,a4paper]{article}

% Pakete
\usepackage[utf8]{inputenc}
\usepackage[T1]{fontenc}
\usepackage[ngerman]{babel}
\usepackage[left=2cm,right=2cm,top=2cm,bottom=2cm]{geometry}
\usepackage{helvet}
\renewcommand{\familydefault}{\sfdefault}
\usepackage{amsmath,amssymb}
\usepackage{graphicx}
\usepackage{tikz}
\usepackage{tcolorbox}
\tcbuselibrary{skins,breakable}
\usepackage{enumitem}
\usepackage{tabularx}
\usepackage{colortbl}
\usepackage{qrcode}
\usepackage{hyperref}
\hypersetup{colorlinks=true,linkcolor=blue,urlcolor=blue}
\usepackage{siunitx}
\sisetup{locale=DE, per-mode=fraction}

% Fachfarbe Physik
\definecolor{physik}{HTML}{2c5aa0}
\definecolor{physiklight}{HTML}{e3f2fd}
\definecolor{loesung}{HTML}{c62828}

% Box-Stile
\tcbset{
    infobox/.style={
        colback=physiklight,
        colframe=physik,
        fonttitle=\bfseries,
        title=#1,
        boxrule=1pt,
        arc=2pt
    },
    warnbox/.style={
        colback=yellow!10,
        colframe=orange!80!black,
        fonttitle=\bfseries,
        title=#1,
        boxrule=1pt,
        arc=2pt
    },
    tipbox/.style={
        colback=green!5,
        colframe=green!50!black,
        fonttitle=\bfseries,
        title=#1,
        boxrule=1pt,
        arc=2pt
    }
}

% Kopfzeile
\usepackage{fancyhdr}
\pagestyle{fancy}
\fancyhf{}
\lhead{\small Physik 12 GK -- Lehrerhinweise}
\chead{\small Photoeffekt Stunde 2}
\rhead{\small 09.01.2026}
\renewcommand{\headrulewidth}{0.4pt}

% Keine Einrückung
\setlength{\parindent}{0pt}
\setlength{\parskip}{6pt}

\begin{document}

% Titel
\begin{center}
{\Large\textbf{Lehrerhinweise: Quantenhypothese \& Gegenfeldmethode}}\\[0.3em]
{\normalsize Stunde 2 (90 min DS) | Fr 09.01.2026 | E = hf, Photogleichung}
\end{center}

\vspace{0.3em}

% ============================================================
% LERNZIELE
% ============================================================

\section*{Feinziele dieser Stunde}

\begin{enumerate}[nosep]
\item SuS \textbf{erklären} die Gegenfeldmethode zur Messung der kinetischen Energie (AFB II)
\item SuS \textbf{berechnen} die Photonenenergie mit $E = h \cdot f$ (AFB II)
\item SuS \textbf{wenden} die Photogleichung $E_{\text{kin}} = hf - W_A$ qualitativ an (AFB II)
\item SuS \textbf{begründen}, warum nur die Frequenz die Grenzspannung bestimmt (AFB II)
\end{enumerate}

\vspace{0.5em}

% ============================================================
% SIMULATION MIT QR-CODE
% ============================================================

\section*{Simulation}

\begin{tcolorbox}[infobox={PhET Photoeffekt -- Gegenfeld-Funktion}]
\begin{minipage}{0.7\textwidth}
\textbf{URL:} \url{https://phet.colorado.edu/de/simulations/photoelectric}

\textbf{Einstellung für diese Stunde:}
\begin{itemize}[nosep]
\item Metall: \textbf{Natrium} (niedrige $W_A$, gut für Berechnungen)
\item Gegenspannung aktivieren (Regler rechts)
\item SuS messen $U_G$ für verschiedene Frequenzen
\end{itemize}

\textbf{Tipp:} Zuerst ohne Gegenspannung zeigen, dann langsam erhöhen bis Strom auf Null geht.
\end{minipage}
\hfill
\begin{minipage}{0.25\textwidth}
\centering
\qrcode[height=2.5cm]{https://phet.colorado.edu/de/simulations/photoelectric}

\vspace{0.2cm}
{\scriptsize PhET Photoeffekt}
\end{minipage}
\end{tcolorbox}

\vspace{0.5em}

% ============================================================
% STUNDENVERLAUF
% ============================================================

\section*{Stundenverlauf (90 min)}

\begin{tabularx}{\textwidth}{|c|l|X|c|c|c|}
\hline
\rowcolor{physiklight}
\textbf{Min} & \textbf{Phase} & \textbf{Aktivität} & \textbf{SF} & \textbf{AB} & \textbf{Folie} \\
\hline
5 & Einstieg & Wiederholung: ,,Was haben wir beobachtet?'' $\rightarrow$ Widerspruch Wellenmodell & PL & -- & F1 \\
\hline
10 & Input & L erklärt Einstein 1905: Licht als Energiepakete, Nobelpreis 1921 & PL & K1 & F2 \\
\hline
10 & Formel & $E = h \cdot f$ einführen, $h$ als Konstante, Einheit eV erklären & PL & K1 & F2 \\
\hline
15 & Übung 1 & SuS berechnen Photonenenergie für verschiedene Farben & EA & A1 & -- \\
\hline
5 & Kontrolle & Ergebnisse vergleichen: Farbspektrum $\leftrightarrow$ Energie & PL & -- & -- \\
\hline
\multicolumn{6}{|c|}{\cellcolor{gray!20}\textit{--- kurze Pause / Übergang ---}} \\
\hline
10 & Gegenfeld & L erklärt Gegenfeldprinzip in Simulation: ,,E-Feld bremst Elektronen'' & PL & K2 & F3 \\
\hline
15 & Messung & SuS messen $U_G$ für 3 Frequenzen in Simulation (Tablets) & PA & A2 & F4 \\
\hline
10 & Ableitung & L leitet Photogleichung ab, zeichnet Einsteinsche Gerade an Tafel & PL & K3 & F5 \\
\hline
5 & Übung 2 & SuS rechnen Aufgabe 3 a--c (Grenzfrequenz, $E_{\text{kin}}$, $U_G$) & EA & A3 & -- \\
\hline
5 & Sicherung & Formeln + Merksätze ins AB übertragen, Vergleich Ergebnisse & EA & K1-3 & F6 \\
\hline
\end{tabularx}

\vspace{0.5em}

\textbf{Legende:} PL = Plenum, EA = Einzelarbeit, PA = Partnerarbeit, K = Kasten, A = Aufgabe, F = Folie (SLIDES-02)

\newpage

% ============================================================
% TAFELBILDER
% ============================================================

\section*{Tafelbilder}

\begin{tcolorbox}[infobox={Tafelbild 1: Einsteins Idee}]
\begin{center}
\textbf{Einsteins Lichtquanten-Hypothese (1905)}
\end{center}

\vspace{0.3em}

Licht = Strom von \textbf{Photonen} (Energiepaketen)

\vspace{0.3em}

$$\boxed{E_{\text{Photon}} = h \cdot f}$$

\vspace{0.3em}

\begin{tabular}{ll}
$h$ & $= \SI{6,63e-34}{J.s} = \SI{4,14e-15}{eV.s}$ (Plancksches Wirkungsquantum) \\
$f$ & = Frequenz des Lichts \\
$E$ & = Energie \textbf{eines} Photons \\
\end{tabular}

\vspace{0.5em}

\textbf{Kernidee: Energie ist gequantelt!}
\begin{itemize}[nosep]
\item Ein Photon = \textbf{ein} Energiepaket = $h \cdot f$
\item \textbf{Unteilbar} -- ganz oder gar nicht (wie eine Münze)
\item Mehr Intensität = mehr Photonen, \textbf{nicht} mehr Energie pro Photon
\end{itemize}

\vspace{0.3em}

\textbf{Beispiel:} Rot ($f = \SI{5e14}{Hz}$): $E = \SI{2,1}{eV}$ \quad Blau ($f = \SI{7,5e14}{Hz}$): $E = \SI{3,1}{eV}$
\end{tcolorbox}

\vspace{0.5em}

\begin{tcolorbox}[infobox={Tafelbild 2: Einsteinsche Gerade ($E_{\text{kin}}$-$f$-Diagramm)}]
\begin{center}
\begin{tikzpicture}[scale=0.85]
% Skalierung: x: 1 Einheit = 10^14 Hz, y: 1 Einheit = 1 eV
% Natrium: W_A = 2,3 eV, f_G = 5,56·10^14 Hz, Steigung = h = 0,414 eV/(10^14 Hz)

% Achsen
\draw[->] (0,0) -- (8.5,0) node[right] {$f$ in $10^{14}$ Hz};
\draw[->] (0,-2.8) -- (0,1.5) node[above] {$E_{\text{kin}}$ in eV};

% Achsenbeschriftung x
\foreach \x in {1,2,3,4,5,6,7,8} {
    \draw (\x,0.08) -- (\x,-0.08);
    \node[below] at (\x,-0.1) {\scriptsize \x};
}
% Achsenbeschriftung y
\foreach \y in {-2,-1,1} {
    \draw (0.08,\y) -- (-0.08,\y);
    \node[left] at (-0.1,\y) {\scriptsize \y};
}
\node[below left] at (0,0) {\scriptsize 0};

% EINE durchgehende Gerade: y = 0.414·x - 2.3
% Extrapolation (gestrichelt) von (0,-2.3) bis (5.56,0)
\draw[dashed, thick, gray] (0,-2.3) -- (5.56,0);
% Messgerade (durchgezogen) von (5.56,0) bis (7.5,0.8)
\draw[very thick, physik] (5.56,0) -- (7.5,0.8);

% Messpunkte (berechnet: E = 0.414·f - 2.3)
\fill[physik] (6.0,0.18) circle (3pt) node[above right, font=\scriptsize] {grün};
\fill[physik] (6.7,0.47) circle (3pt) node[above right, font=\scriptsize] {blau};
\fill[physik] (7.5,0.81) circle (3pt) node[above right, font=\scriptsize] {violett};

% Grenzfrequenz f_G
\fill[red] (5.56,0) circle (3pt);
\draw[dashed, red] (5.56,0) -- (5.56,-0.4);
\node[below, red, font=\small] at (5.56,-0.5) {$f_G = \SI{5,56e14}{Hz}$};

% y-Achsenabschnitt -W_A
\fill[orange] (0,-2.3) circle (3pt);
\node[right, orange, font=\small] at (0.1,-2.3) {$-W_A = \SI{-2,3}{eV}$};

% Steigung einzeichnen
\draw[thick, green!50!black] (6.2,0.26) -- (7.0,0.26) -- (7.0,0.59);
\node[below, green!50!black, font=\scriptsize] at (6.6,0.26) {$\Delta f$};
\node[right, green!50!black, font=\scriptsize] at (7.0,0.42) {$\Delta E$};
\node[right, green!50!black, font=\small] at (7.1,0.7) {Steigung $= h$};
\end{tikzpicture}
\end{center}

\vspace{-0.5em}
$$E_{\text{kin}} = h \cdot f - W_A \quad \text{mit } h = \SI{4,14e-15}{eV.s}$$

\textbf{Interpretation (Natrium):}
\begin{itemize}[nosep]
\item \textbf{Steigung:} $h = \SI{4,14e-15}{eV.s}$ (universell, materialunabhängig!)
\item \textbf{$x$-Achsenabschnitt:} $f_G = W_A / h = \SI{5,56e14}{Hz}$ (Grenzfrequenz)
\item \textbf{$y$-Achsenabschnitt:} $-W_A = \SI{-2,3}{eV}$ (gestrichelt extrapoliert)
\end{itemize}
\end{tcolorbox}

\vspace{1em}

% ============================================================
% DIFFERENZIERUNG
% ============================================================

\section*{Differenzierung}

\begin{tabularx}{\textwidth}{|c|X|X|}
\hline
\rowcolor{physiklight}
\textbf{Niveau} & \textbf{Fokus} & \textbf{Hilfestellung} \\
\hline
Basis & Aufgabe 1 (Einsetzen), Aufgabe 4 a,c,f & Formelkarte: $E = h \cdot f$ mit Beispielrechnung \\
\hline
Standard & Aufgabe 1--4 vollständig & Selbstständig, bei Bedarf Zwischenkontrolle \\
\hline
Erweitert & Zusatzaufgabe (Zink, Wellenlänge) & Zusätzlich: Warum ist $\frac{h}{e}$ die Steigung? \\
\hline
\end{tabularx}

\vspace{0.5em}

\begin{tcolorbox}[tipbox={Tipp für schnelle SuS}]
Wer früher fertig ist: PhET-Simulation mit \textbf{verschiedenen Metallen} wiederholen und $W_A$ aus der Grenzfrequenz berechnen. Erstelle eine Tabelle: Metall | $f_G$ | $W_A$.
\end{tcolorbox}

\vspace{0.5em}

% ============================================================
% HÄUFIGE FEHLER
% ============================================================

\begin{tcolorbox}[warnbox={Häufige Schülerfehler \& Reaktionen}]
\begin{tabularx}{\textwidth}{lX}
\textbf{Fehler} & \textbf{Reaktion} \\
\hline
Einheit eV unklar & ,,1 eV = Energie, die 1 Elektron bei 1 V Spannung aufnimmt'' \\
\hline
$h$ als ,,zu klein'' empfunden & Quantenwelt ist klein $\rightarrow$ kleine Konstante. Aber: Effekte sind messbar! \\
\hline
Gegenfeld-Prinzip abstrakt & Analogie: Ball bergauf rollen $\rightarrow$ verbraucht Energie $\rightarrow$ bleibt stehen \\
\hline
$f \leftrightarrow \lambda$ verwechselt & Merkregel: Hohe Frequenz = kurze Welle = MEHR Energie (,,kurz und heftig'') \\
\hline
$U_G \propto f$ (falsch) & Korrektur: $U_G \propto (f - f_G)$, linear aber mit Offset \\
\hline
\end{tabularx}
\end{tcolorbox}

\newpage

% ============================================================
% LÖSUNGEN
% ============================================================

\section*{Lösungen zum Arbeitsblatt}

\textbf{Aufgabe 1: Photonenenergie berechnen}

\begin{tabular}{|l|c|c|c|}
\hline
\textbf{Lichtart} & \textbf{$f$} & \textbf{Rechnung} & \textbf{$E$} \\
\hline
Rot & $\SI{5,0e14}{Hz}$ & {\color{loesung} $4,14 \cdot 10^{-15} \cdot 5,0 \cdot 10^{14}$} & {\color{loesung} \SI{2,1}{eV}} \\
\hline
Gelb & $\SI{6,0e14}{Hz}$ & {\color{loesung} $4,14 \cdot 10^{-15} \cdot 6,0 \cdot 10^{14}$} & {\color{loesung} \SI{2,5}{eV}} \\
\hline
Blau & $\SI{7,5e14}{Hz}$ & {\color{loesung} $4,14 \cdot 10^{-15} \cdot 7,5 \cdot 10^{14}$} & {\color{loesung} \SI{3,1}{eV}} \\
\hline
UV & $\SI{10,0e14}{Hz}$ & {\color{loesung} $4,14 \cdot 10^{-15} \cdot 10,0 \cdot 10^{14}$} & {\color{loesung} \SI{4,1}{eV}} \\
\hline
\end{tabular}

\vspace{0.5em}

\textbf{Aufgabe 2: Simulation (Werte können leicht variieren)}

\begin{tabular}{|c|c|c|}
\hline
$f$ in $\SI{e14}{Hz}$ & $U_G$ & $E_{\text{kin}}$ \\
\hline
6,0 & {\color{loesung} $\approx$ \SI{0,2}{V}} & {\color{loesung} \SI{0,2}{eV}} \\
\hline
8,0 & {\color{loesung} $\approx$ \SI{1,0}{V}} & {\color{loesung} \SI{1,0}{eV}} \\
\hline
10,0 & {\color{loesung} $\approx$ \SI{1,8}{V}} & {\color{loesung} \SI{1,8}{eV}} \\
\hline
\end{tabular}

\vspace{0.5em}

\textbf{Aufgabe 3: Photogleichung anwenden}

\begin{enumerate}[label=\alph)]
\item {\color{loesung} $f_G = \dfrac{W_A}{h} = \dfrac{\SI{2,3}{eV}}{\SI{4,14e-15}{eV.s}} = \SI{5,56e14}{Hz}$}

\item {\color{loesung} $E_{\text{kin}} = hf - W_A = \SI{4,14}{eV} - \SI{2,3}{eV} = \SI{1,84}{eV} \approx \SI{1,8}{eV}$}

\item {\color{loesung} $U_G = \SI{1,8}{V}$ (numerisch gleich in V und eV)}

\item {\color{loesung} $f_{\text{rot}} = \SI{5,0e14}{Hz} < f_G = \SI{5,6e14}{Hz}$. Photonenenergie (\SI{2,1}{eV}) $<$ Austrittsarbeit (\SI{2,3}{eV}).}
\end{enumerate}

\vspace{0.5em}

\textbf{Aufgabe 4: Konzeptverständnis}

\begin{tabular}{cl}
a) & {\color{loesung} $\boxtimes$ Richtig -- mehr Photonen $\rightarrow$ mehr Elektronen} \\
b) & {\color{loesung} $\square$ Falsch -- Intensität ändert nicht $E$ pro Photon} \\
c) & {\color{loesung} $\boxtimes$ Richtig -- höhere $f$ $\rightarrow$ höhere $E_{\text{kin}}$} \\
d) & {\color{loesung} $\square$ Falsch -- $W_A$ ist Materialeigenschaft} \\
e) & {\color{loesung} $\square$ Falsch -- $U_G \propto (f - f_G)$, nicht $\propto f$} \\
f) & {\color{loesung} $\boxtimes$ Richtig -- bei $f_G$ gerade $E_{\text{kin}} = 0$} \\
\end{tabular}

\vspace{0.5em}

\textbf{Zusatzaufgabe:}

\begin{enumerate}[label=\alph)]
\item {\color{loesung} $f_G = \SI{10,4e14}{Hz}$, $\lambda_G = \dfrac{c}{f_G} = \SI{290}{nm}$ (UV-Bereich)}

\item {\color{loesung} $\lambda_G$ liegt im UV. Sichtbares Licht (400--700 nm) hat zu niedrige $f$ / zu große $\lambda$. Photonen haben $< \SI{4,3}{eV}$, reicht nicht für $W_A$ von Zink.}
\end{enumerate}

\vspace{1em}

% ============================================================
% MATERIAL-CHECKLISTE
% ============================================================

\section*{Material-Checkliste}

\begin{itemize}[nosep]
\item[$\square$] Beamer + Lehrerrechner (PhET-Simulation)
\item[$\square$] Tablets/Laptops für SuS (mind. 1 pro 2 SuS)
\item[$\square$] WLAN-Zugang sicherstellen
\item[$\square$] AB-02 ausgedruckt (1x pro SuS)
\item[$\square$] Tafel/Whiteboard für Diagramm
\item[$\square$] Taschenrechner (falls keine Tablets)
\end{itemize}

\vspace{0.5em}

\hrule
\vspace{0.3em}
\begin{center}
\small Stand: 08.01.2026 | Physik 12 GK
\end{center}

\end{document}
